% !TeX root = ../thesis.tex
% TH-CH04-R1: independent preview compiled; awaiting PDF review, not frozen.
\chapter{多源示教数据构建与质量控制}
\label{chap:multimodal_data}

\section{多模态数据表示与组织方式}
\label{sec:data_representation}

专家轨迹能够完成操作，并不意味着其保存结果已经适合训练\cite{mandlekar2022demonstrations}。模型需要知道某一观测对应哪个动作、哪些量可以作为输入，以及不同记录是否来自相同执行条件。因此，本章在第三章的示教生成基础上，研究状态、动作、视觉和触觉如何组织为可追溯样本，并通过版本与质量标记区分数据用途。这里的多源既包括不同任务和记录路径，也包括独立采集的真实触觉资料；不同来源不因具有相同维度就自动成为同步配对。

数据组织以Episode和Frame为基本单位。Episode表示一次执行及其来源记录，Frame表示该序列中保存的一行观测或转移。Root用于表示按具体规则划分的初态与来源组，不能直接由Episode数量或随机seed数量推导。一个源执行可以派生多个后缀、窗口或导出版本，文件数增加并不意味着独立物理条件增加。对于需要检验泛化的模型，还应根据实际来源确定分组边界，而不是先把所有Frame随机打散再划分数据。

策略可观测状态与完整仿真状态分开保存。当前39维state由受控执行器的13维actuator\_length、13维actuator\_velocity和13维上一位置命令构成，其顺序与机械臂七通道、手部六通道一致。机械臂坐标按角度表达，固定腱坐标按模型长度表达，相应速度单位也不同。该状态不是39个关节位置，不包含完整qpos中的全部自由物体状态；它具有对应真实机器人接口的设计意图，但不因此获得实机信号标定认证。

\begin{table}[htbp]
 \centering\small
 \caption{主要数据字段的维度与实际语义}
 \label{tab:ch04_fields}
 \begin{tabular}{p{0.25\textwidth}p{0.18\textwidth}p{0.43\textwidth}}
 \hline
 字段 & 组织方式 & 语义与单位边界\\
 \hline
 state & float32，$[T,39]$ & 执行器位置、速度及上一命令，非完整关节状态\\
 position\_action & float32，$[T,13]$ & 机械臂关节与手部执行器位置目标\\
 ik\_action & float32，$[T,13]$ & xyz、四元数和手部命令，非实测末端位姿\\
 sensor\_tactile & float32，$[T,712]$ & processed响应，无已认证Pa单位\\
 object\_state & float32，$[T,N,7]$ & 位置与四元数，条目可含固定实体或目标site\\
 qpos、qvel & $[T,D_q]$、$[T,D_v]$ & 完整仿真状态，维度随模型和物体变化\\
 timestamp等 & 每控制行一个时间值 & 仿真时钟及转移结束时间\\
 RGB & 独立视频或图像特征 & 需由时间及控制行索引关联\\
 \hline
 \end{tabular}
\end{table}

表\ref{tab:ch04_fields}中两种动作都为十三维，但语义不同。position\_action保存实际低层位置目标；ik\_action前七维为末端位置与四元数，后六维为手部命令。将IK目标当作关节位置不仅改变单位，还改变学习目标的含义。类似地，目标位姿不能作为实际测量的末端位姿使用，object\_state中的目标site也不应计为独立运动物体。触觉沿用第二章定义的712维processed观测，记录层不重新赋予其真实压力含义。

\begin{figure}[htbp]
 \centering
 \fbox{\begin{minipage}{0.91\textwidth}\centering
 原始执行：source NPZ $+$ Manifest\\[4pt]
 $\downarrow$ 字段、时序与来源核对\\[4pt]
 规范modalities $+$ Episode元数据 $+$ 验证记录\\[4pt]
 $\downarrow$ 按具体导出合同映射\\[4pt]
 LeRobot数值特征、图像或视频、Episode索引\\[4pt]
 真实触觉独立记录另行登记，不自动形成Real--Sim物理配对
 \end{minipage}}
 \caption{数据载体及其转换关系示意}
 \label{fig:ch04_carriers}
\end{figure}

原始NPZ保存密集数组，Manifest记录来源和配置，规范modalities提供统一字段，LeRobot则提供后续训练使用的特征与索引组织，关系见图\ref{fig:ch04_carriers}。这些载体并非互相独立的数据集，每次转换都需要保留父来源。格式版本还应与语义合同分开：LeRobot的codebase\_version不等于五任务在线记录版本。真实触觉窗口也不具备仿真Episode的全部动作和RGB字段，不能为了统一存储而假定缺失模态已经存在。

\section{示教数据采集与关联机制}
\label{sec:data_recording_alignment}

观测与动作的关联决定模型实际学习的问题。以执行前观测预测即将施加的动作，与给模型提供动作执行后的状态，虽然可以生成相同长度的数组，却具有不同的因果含义。若将当前动作与其造成的post-action结果直接配对，训练过程可能利用执行结果反推动作，而不是根据执行前状态选择控制。因此，采集程序需要明确每一行位于控制周期的哪个位置，不能仅依赖行号和时间递增判断关联正确。

设$u_i$为第$i$次施加的动作，$o_i^{-}$与$o_i^{+}$分别为执行前后观测。在线记录采用执行前输入；对于只保存执行后观测的历史序列，在物理连续、无隐藏重置的条件下，可使用前一行建立对应：
\begin{equation}
 \mathcal D_{\mathrm{online}}=\{(o_i^{-},u_i)\},\qquad
 \mathcal D_{\mathrm{shift}}=\{(o_{i-1}^{+},u_i)\}.
 \label{eq:ch04_causal_pairing}
\end{equation}
式\ref{eq:ch04_causal_pairing}不是对所有旧数据都成立的自动修复规则。若中间发生状态恢复、未保存的试探步骤或时间缺口，前一条记录未必就是当前动作的真实起点，需要结合实际执行路径检查。移位还要求源动作语义正确，不能修复把IK输入当成关节命令的问题。

在线记录适配器在物理推进之前采集仿真时钟、策略状态、触觉和对象姿态，推进结束后补充next\_timestamp与实际位置命令。它同时适配普通step和低层step\_physics路径，只将任务记录器最终接受的步骤提交到正式序列，避免把规划器内部探测过程混入示教。低层代码可能先写控制缓冲区再调用物理推进，因此提供prepare\_recorded\_control接口，在覆盖命令之前保存执行前状态。否则，state中的上一命令可能已被当前命令替换，破坏其原有定义。

\begin{figure}[htbp]
 \centering
 \fbox{\begin{minipage}{0.91\textwidth}\centering
 在线：读取$o_i^{-}$ $\rightarrow$ 施加$u_i$ $\rightarrow$ 得到$o_i^{+}$\\[5pt]
 训练配对：$o_i^{-}$ $+$ $u_i$\\[5pt]
 旧记录：保存$o_{i-1}^{+}$ $\rightarrow$ 当前动作$u_i$ $\rightarrow$ 保存$o_i^{+}$\\[5pt]
 合法连续时可配对：$o_{i-1}^{+}$ $+$ $u_i$\\[5pt]
 不直接配对当前执行后结果$o_i^{+}$与动作$u_i$
 \end{minipage}}
 \caption{执行前采集与历史执行后记录的索引关系}
 \label{fig:ch04_timing}
\end{figure}

触觉处理还具有内部时间状态。同一物理状态被多次读取时，如果每次都推进低通滤波或重新采样噪声，新增观测操作本身就可能改变保存值。当前适配器以仿真时间和原始读数内容作为缓存键，同一键只处理一次并返回副本，使额外采集不重复推动信号处理。该机制解释了为什么记录层不仅需要调用传感器接口，还需管理读取时机；它是当前代码的行为，不能回填为所有历史记录均具有的属性。

状态、动作和触觉按控制周期保存。$t_i$对应第$i$条被接受记录在执行前采集的timestamp，$t_i^{+}$对应该控制步推进结束后保存的next\_timestamp；$t_{i+1}$则是下一条被接受记录的执行前时刻。当前在线合同检查
\begin{equation}
 t_{i+1}-t_i\simeq\Delta t_c,\qquad
 t_i^{+}-t_i\simeq\Delta t_c,\qquad \Delta t_c=0.05~\mathrm{s},
 \label{eq:ch04_time_interval}
\end{equation}
其中近似相等表示按实现容差比较浮点时间。仅在相邻控制步连续执行且均被记录、期间没有额外物理推进或重置时，$t_{i+1}$才与$t_i^{+}$一致。未接受步骤、跳步或异常中断可能使两者不再衔接，不能将此关系推广到存在记录缺口的序列；在线检查会拒绝不满足连续时间合同的数据。视频PTS可使用相对帧时间，仿真时间则保留绝对时钟，两者需要由首帧偏移和采样索引连接。时间单调只能排除部分结构错误，不能证明每张图像确实来自对应物理状态。

\begin{figure}[htbp]
 \centering
 \fbox{\begin{minipage}{0.91\textwidth}\centering
 20~Hz控制行：0，1，2，3，4，5，6，7，8，$\ldots$\\[5pt]
 状态与触觉：每个被接受的控制行记录\\[5pt]
 默认20~Hz RGB：0，1，2，3，4，5，6，7，8，$\ldots$\\[5pt]
 5~Hz RGB示例：0，4，8，$\ldots$；保留对应控制行索引\\[5pt]
 图像降频不删除中间控制行，不填造缺少的图像
 \end{minipage}}
 \caption{控制记录与RGB采样频率分离的合同示意}
 \label{fig:ch04_multirate}
\end{figure}

RGB具有独立采样配置。默认合同为1920$\times$1080、20~Hz，速度配置可降至640$\times$360、5~Hz，状态和控制仍为20~Hz。图\ref{fig:ch04_multirate}说明，低频图像需要保存其控制行索引，而不是删除两张图像之间的动作。程序按显式合同检查采样位置，不能把默认视频频率当作全部服务器数据的实测参数。

\begin{table}[htbp]
 \centering\small
 \caption{不同记录路径的时序与版本边界}
 \label{tab:ch04_recording_versions}
 \begin{tabular}{p{0.21\textwidth}p{0.34\textwidth}p{0.31\textwidth}}
 \hline
 路径 & 关联方式 & 使用限制\\
 \hline
 旧生成记录 & 执行后观测及动作列 & 须核对动作语义和索引偏移\\
 修复版LeRobot导出 & 可执行ctrl与前一源观测 & 仅对实际重导出且核验的数据成立\\
 在线记录 & 执行前采集、执行后补实际控制 & 需检查预备控制与接受步骤是否对应\\
 v4精确辅助合同 & 增加完整积分状态和float64控制 & 支持回放，不替代观测关联审查\\
 \hline
 \end{tabular}
\end{table}

Episode终止与异常中断同样影响可保存范围。在线适配器只有在完成对应物理步骤并被记录器接受后才提交行，视频编码异常需向上层传播，不能继续登记缺失模态。表\ref{tab:ch04_recording_versions}中的版本区别因此不仅是字段增减，还涉及生成顺序和验证范围。修复代码的存在、数据已重新导出以及数据通过重新审计是三个不同状态，旧文件不能仅因程序更新而自动获得新合同标记。

\section{数据格式转换与统一存储}
\label{sec:data_conversion}

统一存储的目标是使不同执行来源具有可识别的字段和验证信息，而不是把所有十三维数组都重命名为action。规范数据包从源position\_action映射训练action，从sensor\_tactile映射训练tactile，保留执行前state、时间与对象信息。源NPZ和Manifest仍随包保存，使转换结果可以回到原记录检查。LeRobot导出则进一步按其特征定义组织数值和图像数据；已经存在规范NPZ包，并不意味着对应LeRobot版本也已生成。

历史导出修复首先处理动作来源。Strict回放所需低层控制来自ctrl，而非旧记录中的笛卡尔action。导出器检查控制序列与源ctrl的尾部对应关系，再选取前一条执行后观测作为相应执行前输入。重新执行导出时，程序直接在step之前读状态和图像，之后才施加动作；非重新执行路径则恢复源观测所在的仿真状态，并使用对应触觉。两条路径需要分别记录，不将恢复后渲染与新物理执行混称。

P31时序修复对等长数组执行观测截去末行、动作截去首行的操作，并保存偏移说明和被丢弃的首动作索引。它改变了配对关系，但不自动把IK动作变为关节命令。因此，动作语义修复和索引修复必须独立检查；仅通过shape、finite或调用名为causal的函数，均不能说明训练目标已经正确。对于字段不完整或来源不明的文件，应保留旧版本身份，而不是补零后强行进入统一格式。

转换一致性以同一字段和同一索引下的数值差评价。设源字段按映射$\pi$取值，转换后字段为$\hat{\boldsymbol x}_i$，使用
\begin{equation}
 E_{\max}=\max_{i,j}\left|\hat{x}_{i,j}-x_{\pi(i),j}\right|
 \label{eq:ch04_mapping_error}
\end{equation}
检查指定映射。误差为零只对所比较字段成立，不能扩大为未检查图像、完整状态或物理回放一致。若变换包含量化或重采样，还需按实际合同设置误差解释，而不能一律要求任意格式变换逐比特相同。

指定Sugar Box源共有477行，修复版LeRobot导出340行。按零基索引，导出动作对应源ctrl第137至476行，状态命令和触觉对应前一源行，即第136至475行。本次材料核验重新比较这340行，动作、state末十三维命令以及对应源tactile的最大绝对差均为0。该结果为单一明确来源提供了可复核映射，不代表全部39维状态、RGB像素或其他导出目录均已认证。

v4规范包还保存float64完整积分状态及实际控制，供后续回放使用，这些辅助字段与策略输入分开。Manifest与Episode元数据记录源文件、转换合同及SHA256，文件Hash用于确认文件是否改变，来源说明用于解释其语义。两者缺一不可：相同字段名称或相同文件Hash并不能替代采集与索引规则，统一存储需要的是可追溯合同，而不是只有统一目录形式。

\section{轨迹质量审计与数据清洗}
\label{sec:data_quality_control}

数据质量不能由单一成功标签决定。任务执行成功说明机器人在给定条件下完成目标，Strict通过说明对应轨迹满足特定稳定性约束；两者均未直接检查训练动作的含义或观测时序。反过来，数组finite、shape合法也不说明轨迹完成任务。因此，本章将质量控制分为执行与验证状态、结构数值、时间关联、来源版本、重复隔离以及训练角色六个层次，使每一层回答明确的问题，而不将全部检查折叠成一个模糊的有效标记。

\begin{table}[htbp]
 \centering\small
 \caption{分层数据质量检查及其判定范围}
 \label{tab:ch04_quality_gates}
 \begin{tabular}{p{0.18\textwidth}p{0.36\textwidth}p{0.32\textwidth}}
 \hline
 层次 & 检查对象 & 不能单独推出的结论\\
 \hline
 执行与验证 & 任务结果、Strict配置及状态 & 不证明动作导出和关联正确\\
 结构与数值 & shape、dtype、finite、触觉非负 & 不证明物理合理或任务成功\\
 时间与动作 & 时间间隔、上一命令、物理状态连续 & 不认证全部视频像素同步\\
 来源与版本 & Manifest、源Hash、模型和合同 & 不自动证明跨版本物理等价\\
 重复与隔离 & lineage、content及队列状态 & 不保证独立Root或无近似重复\\
 训练角色 & 标签、组划分及来源用途 & 不由格式合法自动获得训练准入\\
 \hline
 \end{tabular}
\end{table}

结构检查约束state、position\_action、ik\_action和sensor\_tactile的形状与float32类型，检查对象姿态、数值有限性、触觉非负及时间增量。对于已声明版本的阶段编码，还需保证阶段数组长度一致、编码不重复且观测到的编码具有合法含义。字段缺失不能默认为零，尤其不能把未记录触觉伪装成全零触觉。时间和命令连续性检查则需要相应字段及开关实际启用，不能从函数返回的版本名称直接推断所有检查都已执行。

在线状态中的上一命令可以与前一行实际动作比较，保存的执行前qpos/qvel也可与上一行执行后状态比较。这些关系可检出部分记录错位或隐藏状态跳变，但不能替代完整物理回放。旧PickPlace案例中，state后十三维与前一position\_action比较，有373个相邻转移超过$10^{-6}$容差，最大差约0.1。这一发现说明该记录不满足当前命令连续性检查，并不能在没有进一步定位时断言其全部模态反因果或指定唯一故障原因。

旧Stack案例则没有pre\_qpos和pre\_qvel，无法完成当前要求的物理连续性核验。它与旧PickPlace还缺少规范modalities文件，不能描述为所有旧记录都已经完成统一导出。这里的处理是保留问题和版本身份，而不是在本轮删除、补写或重建数据。审计发现的缺陷应与影响范围一起保存，否则后续程序可能仅因旧文件能够读取就重新把它当作已修复数据。

\begin{figure}[htbp]
 \centering
 \fbox{\begin{minipage}{0.91\textwidth}\centering
 执行结果与验证状态 $\rightarrow$ 结构、数值及阶段合同\\[4pt]
 $\downarrow$\\[4pt]
 时间关联与动作语义 $\rightarrow$ 来源、版本与文件Hash\\[4pt]
 $\downarrow$\\[4pt]
 内容重复与用途检查 $\rightarrow$ VALIDATING后登记\\[4pt]
 $\downarrow$\\[4pt]
 登记：\mbox{EXPORTED} / \mbox{DUPLICATE}；异常隔离：\mbox{QUARANTINED}\\[4pt]
 训练准入仍需对应标签、分组与研究用途
 \end{minipage}}
 \caption{分层质量审查与数据准入关系}
 \label{fig:ch04_admission}
\end{figure}

来源检查将规范数组与原始在线记录逐项对应，并核对登记文件Hash、源Manifest、模型证明及独立回放记录之间的绑定。完整积分状态和float64控制可以提供更充分的回放起点与命令，但辅助数组存在并不意味着实际回放已经使用它们，仍需相应证明。图\ref{fig:ch04_admission}中的各层因而既检查数据，也检查执行了哪些验证；代码具有某项能力和某个数据包通过该项验证需要分开记录。

队列通过任务状态、租约和事务组织提交。worker在有效lease下提交结果，登记要求进入\mbox{VALIDATING}且外部verification明确通过。登记时，来源或内容重复者进入\mbox{DUPLICATE}，无重复者插入Episode表并进入\mbox{EXPORTED}；\mbox{QUARANTINED}属于失败或重试耗尽后的隔离路径，并非成功登记的另一出口。过期租约、重试上限和隔离机制用于处理异常，不能让旧worker继续登记已失去所有权的任务。文件写入、验证与数据库登记是不同步骤，异常恢复需要对照已有文件和状态，而不能仅依据进程是否存在判断数据已完整保存。上述机制是实现事实，不代表任意磁盘断连之后都已经证明无损恢复。

重复检测同时使用来源和内容标识。对于待登记记录$e$及已登记集合$\mathcal A$，当前逻辑可表示为
\begin{equation}
 \operatorname{dup}(e)=\boldsymbol 1\!\left[
 \exists e'\in\mathcal A:\ L(e)=L(e')\ \lor\ C(e)=C(e')\right],
 \label{eq:ch04_duplicate}
\end{equation}
其中，$L$与$C$分别为lineage和content标识。相同来源或相同语义内容任一成立，即进入重复处理，而不是要求两个Hash同时相同。lineage绑定任务、资产与源文件，content则根据规范数组和解码RGB摘要生成；时间使用相对首行的表示，精确回放辅助字段不参与内容摘要，以减少仅改变时钟或多存辅助信息造成的重复计账。

Hash只能识别定义范围内的相等，不能证明两个不同Hash具有不同物理初态，也不能排除近似轨迹或相同模板派生的样本。训练划分还需要按来源组管理，同一Root的窗口、后缀和重导出版本不能因文件不同而分散到训练与验证两侧。对于缺少完整lineage的历史文件，应把独立性标记为未知，而不是用新的随机划分掩盖来源缺口。

数据清洗也不应等同于删除所有失败轨迹。合法失败仍可用于失败诊断或具有明确标签的其他研究，但不能直接当作成功专家标签；是否计入生产时长应遵循对应冻结规则。历史版本保留来源和合同标记，既能支持修复效果比较，也能避免旧性能结论继续建立在错误语义上。表\ref{tab:ch04_quality_gates}所示分层审查因此强调按用途准入：任务成功、Strict通过、格式有效、时序正确与可用于具体训练之间存在依赖，却不能互相替代。

\section{数据集统计与质量验证}
\label{sec:dataset_quality_validation}

本节以明确的数据范围检验上述合同和审查机制。材料核验直接读取三组LeRobot目录的全部选定数值Parquet行，以及十一条指定源NPZ，后者包括五条旧在线记录、五个v4开发包和一个Sink精度对照源。数字均绑定对应目录与版本，不以历史summary替代当前读取结果。对于视频解码、物理回放或真实触觉原始全库，本轮没有重新执行相应验证，其已有合同和成绩分别标明来源层级。

三组LeRobot库存见表\ref{tab:ch04_lerobot_inventory}。实际episode\_index数量和Parquet行数与info中的Episode、Frame声明一致，state、action和tactile分别符合39、13、712维，所读数值finite，触觉非负。这个结果说明指定数值文件可按所声明结构读取，但不证明早期版本的动作语义和时序正确。两个旧目录仍保留其历史合同问题，不能因为当前程序已经修复，就把旧数据和模型性能重新解释为合规基线。

\begin{table}[htbp]
 \centering\small
 \caption{三组LeRobot数据的实际数值库存}
 \label{tab:ch04_lerobot_inventory}
 \begin{tabular}{p{0.30\textwidth}rrp{0.27\textwidth}}
 \hline
 数据版本 & Episode & Frame & 使用范围\\
 \hline
 早期五条导出 & 5 & 1550 & 历史格式及合同问题对照\\
 早期formal30导出 & 30 & 14310 & 保留原版本，不自动认证语义\\
 Sugar单物体修复版 & 1 & 340 & 指定源数值映射核验\\
 \hline
 \end{tabular}
\end{table}

三组数据的触觉非零比例分别约为12.7297\%、9.4029\%和11.0972\%，这是对各目录所选全部数值行的重新计算，不是任务成功率或物理准确率。task\_index分别对应两个、三个和一个lift:物体标签，实际都属于Lift操作，不能解释为覆盖两种、三种和一种物理任务。三个info均仅声明train范围，未提供本轮可认证的独立验证和测试划分。Episode数量因此不能直接替代独立Root数量，也不能将三个版本相加为36个独立采集主体。

修复版的一致性检查进一步回到原始来源。340行导出action与源ctrl、state命令区与前一源ctrl、触觉与前一源tactile均得到最大绝对差0，支持该索引映射在指定数值字段上成立。它不意味着所有39维状态都已单独重算，也不包括RGB像素逐帧检查。对三组数据的完整Episode数值摘要比较未发现完全相同的摘要，但这一检查既不含图像，也不排除截断后缀、轻微差异或同源版本重复，不能据此宣布无泄漏。

\begin{table}[htbp]
 \centering\small
 \caption{旧在线记录与v4开发包的质量对照}
 \label{tab:ch04_version_quality}
 \begin{tabular}{p{0.23\textwidth}p{0.34\textwidth}p{0.29\textwidth}}
 \hline
 对象 & 本次材料核验结果 & 限制\\
 \hline
 旧PickPlace & 373个转移命令关系超容差，最大约0.1 & 不满足当前连续性检查，未修复\\
 旧Stack & 缺pre\_qpos、pre\_qvel及modalities & 物理连续性UNKNOWN\\
 五任务v4开发包 & 2727行，字段映射及登记文件Hash匹配 & 开发准入，非正式生产验收\\
 v4独立性标识 & 5个lineage、4个asset\_hash、5个content值 & 不等于认证独立Root数\\
 \hline
 \end{tabular}
\end{table}

旧五任务在线样本共2665行，表\ref{tab:ch04_version_quality}中的PickPlace与Stack缺陷说明，同一历史批次也不能只依据某个样本通过就整体升级。相对地，所选五个v4包覆盖Lift、PickPlace、Stack、Push和NutAssembly，共2727个控制行。其命令及所存qpos/qvel连续性检查通过，规范action、触觉、state与源映射逐项一致，登记文件Hash全部匹配且无登记文件缺失。这里确认的是保存数据及文件关系，不能替代新的物理执行。

五个v4包有五个lineage、四个登记asset\_hash及五个登记content值，独立Root数仍为UNKNOWN。包中admission明确为开发验证，而非正式生产准入。五个视频文件都存在，只能作为文件覆盖证据；本轮没有对全部RGB解码及逐帧同步重新认证，也没有独立重算包含解码视频摘要的语义content hash。历史视频验证记录可以随版本引用，却不应升级为本轮完成的视频验收。

历史Sink对照提供了回放合同影响的另一实例。同一380帧源在旧float32路径下回放失败，最大qpos分量偏差为1.628897；保存完整积分状态并使用float64控制的路径回放通过，记录偏差为0，见表\ref{tab:ch04_sink_precision}及图\ref{fig:ch04_sink_precision}。这些数值来自已保存回放报告，不是本轮新实验；本轮只读取源数组并确认辅助字段和数据类型。qpos包含不同类型的坐标，最大分量偏差也不能统一解释为某一长度误差。

\begin{table}[htbp]
 \centering\small
 \caption{Sink历史回放状态与精度对照}
 \label{tab:ch04_sink_precision}
 \begin{tabular}{p{0.46\textwidth}rp{0.17\textwidth}}
 \hline
 历史回放路径 & 最大qpos分量偏差 & 历史结果\\
 \hline
 旧float32状态及控制路径 & 1.628897 & 失败\\
 完整积分状态与float64控制路径 & 0 & 通过\\
 \hline
 \end{tabular}
\end{table}

\begin{figure}[htbp]
 \centering
 \includegraphics[width=0.78\textwidth]{figures/ch04/sink_historical_precision.pdf}
 \caption{由已有报告绘制的Sink回放偏差对照（历史记录，非本轮新实验）}
 \label{fig:ch04_sink_precision}
\end{figure}

两种设置同时改变完整状态保存与控制精度，因此该对照支持原状态/精度合同不足的解释，但没有分离两个因素的独立贡献。偏差0也只指报告中对应保存状态的比较方式，不能推广为任意任务、模型版本和浮点后端均逐位复现。保存额外状态能够增加回放证据，却不产生新的物理多样性；同一条件的旧隔离副本和新通过副本不能重复累计为新增时长。

真实触觉与Natural Sim的冻结训练合同需另行登记，不混入上述Episode与Frame库存。已有P44合同声明Real Train为5153窗口、32 Session，Real Validation为780窗口、4 Session；Natural Train为4049窗口、50 Root，Validation为905窗口、10 Root。这些是历史合同声明，本轮未重新认证原始Real全库。合同还保留Real历史不等间隔帧与Natural来源20~Hz连续采样的区别，因此不能把四帧窗口数量解释为同样的物理持续时间或相同数量的独立接触过程。

Real记录与仿真示教没有已认证同步物理配对，物理通道映射也仍有限制。它们可以按来源组织，并用于后续明确监督条件的跨域研究，但不能作为已经同步的视觉、动作与真实触觉联合数据集。Validation名称本身也不消除历史开发暴露，训练角色和独立性需要结合实际使用记录解释，而不是仅凭文件名判断。

本节统计不覆盖服务器批次2的实时数据库和完整磁盘。当前全量有效小时、独立Root及全局泄漏状况均不能由这几组样本外推。重复读取扩容索引反映的是采样使用量，不是新增物理仿真时长。现有证据支持版本化合同能够发现问题、指定转换和开发包满足所列检查；它尚不支持全部历史数据均可训练、全任务模态完全同步或正式生产指标已通过的结论。

\section{本章小结}
\label{sec:data_quality_summary}

本章从训练样本需要具有明确输入、动作和来源关系出发，建立了多模态数据表示、采集关联、规范转换及分层审查的论证。39维策略状态与完整仿真状态分开，两种十三维动作按语义区分，触觉和RGB按各自采样合同关联。执行前记录、控制缓冲预备和触觉缓存共同说明，采集时机本身是数据定义的一部分，不能由格式和维数代替。

已有数据核验确认了指定Sugar导出字段与源记录的一致性，也发现旧PickPlace命令连续性和旧Stack字段缺失问题。五任务v4开发包及Sink历史回放对照提供了版本化质量控制的有限依据，但不构成服务器全量验收、独立Root认证或全局无泄漏证明。保留来源、Hash、验证层级和失败记录，使这些局限能够被明确识别，而非通过统一文件名掩盖。

在上述数据条件下，后续研究需要进一步区分真实触觉与仿真触觉的分布及监督关系。下一章将基于已保存的真实统计信息和冻结触觉子集讨论跨域重建，不将规范化存储等同于两域已经具有同步物理对应。
